% ID: SEL-D0084-OA-E01
% Tanda S084.
\textbf{a)} Al transponer $F$, las filas pasan a corresponder a los artículos. Multiplicamos filas por columnas en cada producto:
\[
 F^tG=\boxed{\begin{pmatrix}
 1400&1800&1100\\1900&2450&1500\\1040&1340&820
 \end{pmatrix}},
 \qquad
 FG^t=\boxed{\begin{pmatrix}2200&1390\\3900&2470\end{pmatrix}}.
\]
Por ejemplo, $(F^tG)_{11}=100\cdot6+200\cdot4=1400$, mientras que $(FG^t)_{11}=100\cdot6+150\cdot8+80\cdot5=2200$.

\textbf{b)} En $F^tG$, las entradas de la diagonal multiplican las cantidades de cada artículo por sus propias ganancias unitarias en ambos formatos. Por tanto, las ganancias de A, B y C son 1400, 2450 y 820 euros, respectivamente.

\textbf{c)} En $FG^t$, las entradas diagonales suman las ganancias de todos los artículos en un mismo formato. El formato grande aporta 2200 euros y el normal 2470 euros. La ganancia total es $\boxed{4670}$ euros. Se comprueba también sumando $1400+2450+820=4670$.
